\chapter{Training data and the training loop: gradients, AdamW, warmup and cosine schedules}\label{ch:trainloop}
\begingroup\emergencystretch=3em % long unbreakable source tags; local to this chapter

\begin{onething}
Training repeats one small move millions of times: take a batch of patches, compute the loss
of \cref{ch:losses}, compute its gradient with respect to every parameter, and let the
optimiser (here Adam) step each parameter a little against its gradient, with a step size
that a schedule (here cosine annealing) shrinks over the run. Everything else in the loop,
from shards to checkpoints, exists to keep that move fed, honest and resumable.
\end{onething}

\Cref{ch:losses} ended with a number, the loss, and its derivative with respect to the
reconstruction. Two gaps remain between that and a trained model. First, the network has about
5.5 million parameters \tagM{research/model_break_down/05_hyperspectral_segformer_mae.md}, not
one reconstruction; you need the derivative with respect to each of them and a rule for what
to do with it. Second, the loss is computed on one batch of patches, and the satellite archive
holds millions; you need a pipeline that turns scenes into a steady stream of batches, and a
loop that consumes it. This chapter builds both: the mathematics of the update rule first
(\cref{sec:trainloop:gd,sec:trainloop:momentum,sec:trainloop:adam,sec:trainloop:decay,sec:trainloop:schedule,sec:trainloop:clip,sec:trainloop:precision}),
then the data (\cref{sec:trainloop:data}), then the loop that ties them together
(\cref{sec:trainloop:loop}).

The obvious approach, ``try a change to each parameter and keep it if the loss went down'',
fails on arithmetic alone: one trial per parameter is 5.5 million forward passes for a single
improvement. The gradient gives you the best local direction for all 5.5 million parameters at
the cost of about one extra pass.

\begin{analogy}
Training is walking down a fog-bound hillside to the valley floor. You cannot see the valley;
you can feel the slope under your feet (the gradient). Plain gradient descent takes a step
straight downhill each time. Momentum is a heavy ball that keeps rolling through small bumps.
Adam is a walker who shortens their stride on ground that has been steep and jagged and
lengthens it on ground that has been gentle and consistent. Warmup is taking small steps while
your eyes adjust to the fog; cosine annealing is slowing down as you near the bottom so you
stop in the valley instead of striding across it.
\end{analogy}

\section{Vocabulary: parameter, gradient, batch, step, epoch}\label{sec:trainloop:vocab}

\begin{definition}[Gradient]\index{gradient}
For a loss $L$ that depends on parameters $\vect{w} = (w_1, \dots, w_P)$, the \emph{gradient}
is the vector of partial derivatives,
\[
\nabla L = \nabla_{\vect{w}} L = \left(\frac{\partial L}{\partial w_1}, \dots, \frac{\partial L}{\partial w_P}\right).
\]
Entry $j$ says how fast the loss changes when $w_j$ alone is nudged up. The gradient points in
the direction of steepest increase, so $-\nabla L$ points steepest downhill.
\end{definition}

\begin{definition}[Batch, step, epoch]\index{batch}\index{step}\index{epoch}
A \emph{batch} (or \emph{mini-batch}) is the group of patches processed together in one forward
pass; its size $B$ is \code{batch_size}. A \emph{step} (or \emph{iteration}) is one update of
the parameters. In most textbooks an \emph{epoch} is one full pass over the training set; in
this project an epoch is a fixed \emph{sample budget} instead (\cref{sec:trainloop:epoch}),
the interval at which validation, the learning-rate schedule, the SAM ramp and checkpointing
advance.
\end{definition}

The loss over a batch is the mean of the per-patch losses, and the gradient of a mean is the
mean of the gradients (differentiation is linear). So the batch gradient is an
\emph{estimate} of the gradient you would get from the whole training set: noisy, but cheap.
Training with such estimates is \emph{stochastic gradient descent}\index{stochastic gradient
descent}. A larger batch gives a less noisy estimate; it does not give a different
expected direction.

\paragraph{Backpropagation, in one paragraph.} \Cref{ch:losses} computed $\partial L/\partial\xhat$.
The reconstruction $\xhat$ is the output of the decoder, which is a function of its weights and
of the encoder's output, and so on back to the compressor. The chain rule
(\cref{sec:losses:chain}) carries the derivative backwards layer by layer: each layer receives
$\partial L/\partial(\text{its output})$, and computes from it $\partial L/\partial(\text{its
weights})$ and $\partial L/\partial(\text{its input})$, which it passes on. PyTorch does this
automatically when you call \code{loss.backward()}; it stores each parameter's derivative in the
parameter's \code{.grad} field, \emph{adding} to whatever is already there. That adding is
why the loop calls \code{optimizer.zero_grad()} every step, and it is what gradient
accumulation exploits (\cref{sec:trainloop:accum}).

\section{Gradient descent, by hand}\label{sec:trainloop:gd}

\begin{definition}[Gradient descent]\index{gradient descent}\index{learning rate}
With \emph{learning rate} (step size) $\eta > 0$, each step replaces every parameter by
\begin{equation}\label{eq:trainloop:gd}
\vect{w}_{t+1} = \vect{w}_t - \eta\,\nabla L(\vect{w}_t).
\end{equation}
\end{definition}

To see the rule work, use a toy loss with one parameter whose answer you know: $L(w) = (w -
3)^2$, smallest at $w = 3$. Its derivative, by the chain rule, is $L'(w) = 2(w - 3)$.

\begin{worked}[label={we:trainloop:gd}]{Three steps of gradient descent}
Do this on paper. Start at $w_0 = 0$ (loss $(0 - 3)^2 = 9$), with $\eta = 0.1$.

\textbf{Step 1.}
\begin{align*}
g_0 &= 2(0 - 3) = -6\\
w_1 &= 0 - 0.1 \times (-6) = 0 + 0.6 = 0.6\\
L(w_1) &= (0.6 - 3)^2 = (-2.4)^2 = 5.76
\end{align*}
\textbf{Step 2.}
\begin{align*}
g_1 &= 2(0.6 - 3) = 2 \times (-2.4) = -4.8\\
w_2 &= 0.6 - 0.1 \times (-4.8) = 0.6 + 0.48 = 1.08\\
L(w_2) &= (1.08 - 3)^2 = (-1.92)^2 = 3.6864
\end{align*}
\textbf{Step 3.}
\begin{align*}
g_2 &= 2(1.08 - 3) = 2 \times (-1.92) = -3.84\\
w_3 &= 1.08 - 0.1 \times (-3.84) = 1.08 + 0.384 = 1.464\\
L(w_3) &= (1.464 - 3)^2 = (-1.536)^2 = 2.359296
\end{align*}
The loss falls $9 \to 5.76 \to 3.69 \to 2.36$. The negative gradient pointed right (towards 3)
every time, and the steps shrink as the slope flattens.
\end{worked}

\paragraph{Why the steps shrink, and when they explode.} Substitute the derivative into
\cref{eq:trainloop:gd} and subtract 3 from both sides:
\begin{align*}
w_{t+1} &= w_t - 2\eta(w_t - 3)
  && \text{(the update with $L' = 2(w - 3)$)}\\
w_{t+1} - 3 &= (w_t - 3) - 2\eta(w_t - 3)
  && \text{(subtract 3 from both sides)}\\
w_{t+1} - 3 &= (1 - 2\eta)(w_t - 3).
  && \text{(factor out $w_t - 3$)}
\end{align*}
The distance to the answer is multiplied by $1 - 2\eta$ every step. With $\eta = 0.1$ the factor
is $0.8$: $-3 \to -2.4 \to -1.92 \to -1.536$, exactly the distances above. The method converges
when $\abs{1 - 2\eta} < 1$, that is $0 < \eta < 1$. With $\eta = 1.1$ the factor is $-1.2$: each
step overshoots to the other side and lands further away, $w = 6.6, -1.32, 8.184$, and the
loss grows $12.96, 18.66, 26.87$ \tagC. A learning rate too large does not merely slow
training; it makes it diverge. Real networks have millions of such directions with different
curvatures, and the learning rate must suit the steepest of them.

\section{Momentum}\label{sec:trainloop:momentum}

Plain gradient descent forgets everything at every step. On a long, gently sloping valley it
crawls; on a narrow one it zig-zags across. \emph{Momentum}\index{momentum}
\citep{polyak1964some} keeps a running velocity:
\begin{equation}\label{eq:trainloop:momentum}
\vect{v}_{t} = \beta\,\vect{v}_{t-1} + \nabla L(\vect{w}_t), \qquad
\vect{w}_{t+1} = \vect{w}_t - \eta\,\vect{v}_t, \qquad \vect{v}_{-1} = \vect{0},
\end{equation}
with $\beta$ typically 0.9. (This is PyTorch's \code{SGD(momentum=0.9)} form.) Gradients that
keep pointing the same way add up; gradients that alternate cancel.

\begin{worked}[label={we:trainloop:momentum}]{Three steps with momentum}
Same toy, $w_0 = 0$, $\eta = 0.1$, $\beta = 0.9$.

\textbf{Step 1.} $g = 2(0 - 3) = -6$. $v = 0.9 \times 0 + (-6) = -6$. $w = 0 - 0.1 \times (-6) = 0.6$.

\textbf{Step 2.} $g = 2(0.6 - 3) = -4.8$.
\begin{align*}
v &= 0.9 \times (-6) + (-4.8) = -5.4 - 4.8 = -10.2\\
w &= 0.6 - 0.1 \times (-10.2) = 0.6 + 1.02 = 1.62
\end{align*}
\textbf{Step 3.} $g = 2(1.62 - 3) = -2.76$.
\begin{align*}
v &= 0.9 \times (-10.2) + (-2.76) = -9.18 - 2.76 = -11.94\\
w &= 1.62 - 0.1 \times (-11.94) = 1.62 + 1.194 = 2.814
\end{align*}
After three steps $w = 2.814$ (loss $0.0346$), against plain descent's $1.464$ (loss $2.36$).
PyTorch's \code{SGD} with momentum gives $2.814$ \tagC. The velocity is still large, so step 4
will overshoot past 3; on this toy, momentum oscillates around the answer before settling.
\end{worked}

\section{Adam}\label{sec:trainloop:adam}

\subsection{The problem Adam solves}

Different parameters see gradients of wildly different sizes. A bias in the final layer may
see gradients of order 1; a weight deep in the encoder, of order $10^{-4}$. One learning rate
cannot suit both: large enough for the second makes the first diverge. Adam \citep{kingma2015adam}
gives every parameter its own effective step size, by dividing its gradient by a running
estimate of the gradient's typical magnitude.

\subsection{Exponential moving averages}

\begin{definition}[Exponential moving average]\index{exponential moving average}
An \emph{exponential moving average} (EMA) of a sequence $g_1, g_2, \dots$ with decay $\beta$ is
\[
m_t = \beta\,m_{t-1} + (1 - \beta)\,g_t, \qquad m_0 = 0.
\]
\end{definition}
Unroll it to see what it averages:
\begin{align*}
m_1 &= (1-\beta)g_1\\
m_2 &= \beta(1-\beta)g_1 + (1-\beta)g_2\\
m_t &= (1-\beta)\sum_{k=1}^{t}\beta^{t-k}g_k .
\end{align*}
Each older gradient's weight shrinks by a factor $\beta$ per step: with $\beta = 0.9$ a gradient
ten steps old has weight $0.9^{10} = 0.35$ of the newest; with $\beta = 0.999$ the memory is
about a thousand steps long.

\subsection{The algorithm}

\begin{definition}[Adam]\index{Adam}
For each parameter $w$ separately, with gradient $g_t$ at step $t = 1, 2, \dots$, decays
$\beta_1, \beta_2$, learning rate $\eta$ and a small $\epsilon$:
\begin{align}
m_t &= \beta_1 m_{t-1} + (1 - \beta_1)\,g_t && \text{(EMA of the gradient: first moment)}\label{eq:trainloop:m}\\
v_t &= \beta_2 v_{t-1} + (1 - \beta_2)\,g_t^2 && \text{(EMA of the squared gradient: second moment)}\label{eq:trainloop:v}\\
\hat{m}_t &= \frac{m_t}{1 - \beta_1^t}, \qquad \hat{v}_t = \frac{v_t}{1 - \beta_2^t} && \text{(bias correction)}\label{eq:trainloop:bias}\\
w_t &= w_{t-1} - \eta\,\frac{\hat{m}_t}{\sqrt{\hat{v}_t} + \epsilon}. && \text{(the step)}\label{eq:trainloop:adamstep}
\end{align}
PyTorch's defaults, which this project uses unchanged, are $\beta_1 = 0.9$, $\beta_2 = 0.999$,
$\epsilon = 10^{-8}$ and weight decay 0.
\end{definition}

$\hat{m}_t$ estimates the mean gradient; $\sqrt{\hat{v}_t}$ estimates its root-mean-square size.
Their ratio is near $\pm 1$ when the gradient is consistent and near 0 when it flips sign from
step to step. So each parameter moves about $\eta$ per step when its gradient is steady, whatever
the gradient's scale, and less when it is noisy.

\subsection{Why bias correction}

Both EMAs start at 0, so early on they are too small. Suppose the gradient's distribution does
not change, with mean $\E[g]$. Take expectations of the unrolled form:
\begin{align*}
\E[m_t] &= (1 - \beta_1)\sum_{k=1}^{t}\beta_1^{t-k}\,\E[g]
  && \text{(expectation is linear)}\\
&= (1 - \beta_1)\,\E[g]\,(1 + \beta_1 + \dots + \beta_1^{t-1})
  && \text{(write out the powers)}\\
&= (1 - \beta_1)\,\E[g]\,\frac{1 - \beta_1^t}{1 - \beta_1}
  && \text{(sum of a geometric series)}\\
&= (1 - \beta_1^t)\,\E[g].
  && \text{(cancel $1 - \beta_1$)}
\end{align*}
The EMA is short by the factor $1 - \beta_1^t$; dividing by it, as in \cref{eq:trainloop:bias},
removes the bias. The same argument with $g^2$ gives the factor $1 - \beta_2^t$ for $v_t$.

\begin{worked}[label={we:trainloop:bias}]{The bias-correction factors for t = 1, 2, 3}
With $\beta_1 = 0.9$:
\begin{align*}
t = 1:&\quad 1 - 0.9 = 0.1\\
t = 2:&\quad 1 - 0.9^2 = 1 - 0.81 = 0.19\\
t = 3:&\quad 1 - 0.9^3 = 1 - 0.729 = 0.271
\end{align*}
With $\beta_2 = 0.999$:
\begin{align*}
t = 1:&\quad 1 - 0.999 = 0.001\\
t = 2:&\quad 1 - 0.999^2 = 1 - 0.998001 = 0.001999\\
t = 3:&\quad 1 - 0.999^3 = 1 - 0.997003 = 0.002997
\end{align*}
So at $t = 1$ the raw $m_1$ is divided by 0.1 (multiplied by 10) and the raw $v_1$ by 0.001
(multiplied by 1000). Without correction, the first step would divide a gradient estimate that
is 10 times too small by a size estimate that is $\sqrt{1000} = 31.6$ times too small, and the
step would be about 3 times too large. The second-moment correction matters for much longer:
$1 - 0.999^{t}$ is still only $0.095$ at $t = 100$ and $0.632$ at $t = 1000$ \tagC.
\end{worked}

\begin{worked}[label={we:trainloop:adam}]{Three steps of Adam on the toy loss}
Toy $L(w) = (w - 3)^2$, $w_0 = 0$, $\eta = 0.1$, default $\beta$s; ignore $\epsilon$ (it is
$10^{-8}$).

\textbf{Step $t = 1$.} $g_1 = 2(0 - 3) = -6$.
\begin{align*}
m_1 &= 0.9 \times 0 + 0.1 \times (-6) = -0.6\\
v_1 &= 0.999 \times 0 + 0.001 \times 36 = 0.036\\
\hat{m}_1 &= -0.6 / 0.1 = -6\\
\hat{v}_1 &= 0.036 / 0.001 = 36, \qquad \sqrt{36} = 6\\
w_1 &= 0 - 0.1 \times (-6)/6 = 0 + 0.1 = 0.1
\end{align*}
\textbf{Step $t = 2$.} $g_2 = 2(0.1 - 3) = -5.8$, $g_2^2 = 33.64$.
\begin{align*}
m_2 &= 0.9 \times (-0.6) + 0.1 \times (-5.8) = -0.54 - 0.58 = -1.12\\
v_2 &= 0.999 \times 0.036 + 0.001 \times 33.64 = 0.035964 + 0.03364 = 0.069604\\
\hat{m}_2 &= -1.12 / 0.19 = -5.894737\\
\hat{v}_2 &= 0.069604 / 0.001999 = 34.81941, \qquad \sqrt{34.81941} = 5.900797\\
w_2 &= 0.1 - 0.1 \times (-5.894737)/5.900797 = 0.1 + 0.099897 = 0.199897
\end{align*}
\textbf{Step $t = 3$.} $g_3 = 2(0.199897 - 3) = -5.600205$, $g_3^2 = 31.36230$.
\begin{align*}
m_3 &= 0.9 \times (-1.12) + 0.1 \times (-5.600205) = -1.008 - 0.560021 = -1.568021\\
v_3 &= 0.999 \times 0.069604 + 0.001 \times 31.36230 = 0.069534 + 0.031362 = 0.100897\\
\hat{m}_3 &= -1.568021 / 0.271 = -5.786054\\
\hat{v}_3 &= 0.100897 / 0.002997 = 33.66589, \qquad \sqrt{33.66589} = 5.802231\\
w_3 &= 0.199897 - 0.1 \times (-5.786054)/5.802231 = 0.199897 + 0.099721 = 0.299618
\end{align*}
\code{torch.optim.Adam(lr=0.1)} on the same toy gives $w = 0.100000, 0.199897, 0.299618$ \tagC
(\file{research/textbook/tools/compute_24-trainloop_output.txt}).

Notice the steps: $0.1000, 0.0999, 0.0997$. Each is almost exactly $\eta$, although the
gradient was $-6$. Plain descent with the same $\eta$ took a first step of $0.6$.
\end{worked}

The first Adam step is $\eta\,\hat{m}_1/\sqrt{\hat{v}_1} = \eta\,g_1/\abs{g_1} = \pm\eta$ for
every parameter, whatever the size of its gradient: for $g_1 = 10^{-4}$, $1$ or $10^4$ the step is
$0.1$ \tagC. That is Adam's strength (no parameter is starved) and the reason it needs care at
the start: on step 1, every one of the 5.5 million parameters moves by the full learning rate,
in a direction estimated from a single batch. \Cref{sec:trainloop:warmup} returns to this.

\begin{sourcenote}[: the optimiser actually used]
\file{app/abstract_classes/foundation_trainer.py}, lines 68--70:
\code{self.optimizer = torch.optim.Adam(self.model.parameters(), lr=config.learning_rate)}. Plain
Adam with PyTorch's defaults: $\beta_1 = 0.9$, $\beta_2 = 0.999$, $\epsilon = 10^{-8}$,
\code{weight_decay=0}. The fine-tuning trainer builds the same optimiser over the trainable
parameters only (\file{app/foundation_models/trainers/finetune_trainer.py}, line 254). There is
no AdamW anywhere in \file{app/}, \file{scripts/} or \file{research/} (a grep for
\code{AdamW} finds nothing), and no weight decay. The course (IW28) says so and is right.
\end{sourcenote}

\section{Weight decay, L2 regularisation and AdamW}\label{sec:trainloop:decay}

The project does not use weight decay, but the chapter title promises AdamW, and the
distinction is one of the most common confusions in practice, so it is worth deriving.

\begin{definition}[Weight decay]\index{weight decay}
\emph{Weight decay} with coefficient $\lambda_w$ shrinks every weight towards zero a little at
each step: $w \leftarrow w - \eta\lambda_w w$. It discourages large weights, a form of
\emph{regularisation}\index{regularisation} (a preference for simpler models that is not part of
fitting the data).
\end{definition}

\paragraph{With plain gradient descent, weight decay equals an L2 penalty.} Add $\frac{\lambda_w}{2}\norm{\vect{w}}^2$
to the loss. Its gradient is $\lambda_w\vect{w}$, so the update becomes
\[
\vect{w} \leftarrow \vect{w} - \eta\big(\nabla L + \lambda_w\vect{w}\big)
= \vect{w} - \eta\nabla L - \eta\lambda_w\vect{w},
\]
gradient descent followed by weight decay. The two ideas are the same thing here.

\paragraph{With Adam they are not.} If you add $\lambda_w\vect{w}$ to the gradient (``L2 in
Adam''), the decay term enters $m_t$ and $v_t$ and is then divided by $\sqrt{\hat{v}_t}$ like
everything else. A parameter whose gradients have been large gets a large $\sqrt{\hat{v}_t}$, and
therefore \emph{less} decay; a parameter with small gradients gets \emph{more}. AdamW
\citep{loshchilov2019decoupled} \emph{decouples} the decay: it applies Adam to the loss gradient
alone and subtracts $\eta\lambda_w w$ separately,
\begin{equation}\label{eq:trainloop:adamw}
w_t = w_{t-1} - \eta\left(\frac{\hat{m}_t}{\sqrt{\hat{v}_t} + \epsilon} + \lambda_w\,w_{t-1}\right).
\end{equation}

\begin{worked}[label={we:trainloop:adamw}]{The same decay, two different shrinks}\index{AdamW}
Two weights, both $w = 1.0$, with $\lambda_w = 0.01$, $\eta = 10^{-3}$. Weight A has had large
gradients, $\sqrt{\hat{v}} = 10$; weight B small ones, $\sqrt{\hat{v}} = 0.01$. Look only at the
decay part of one step (take the loss gradient to be zero at this moment, and treat the
decay term as passing straight through the normaliser, the steady-state picture).

\textbf{L2 in Adam:} the decay $\lambda_w w = 0.01$ is divided by $\sqrt{\hat{v}}$.
\begin{align*}
\text{A: } & 10^{-3} \times 0.01 / 10 = 10^{-6}\\
\text{B: } & 10^{-3} \times 0.01 / 0.01 = 10^{-3}
\end{align*}
\textbf{AdamW:} the decay is not normalised.
\begin{align*}
\text{A: } & 10^{-3} \times 0.01 \times 1.0 = 10^{-5}\\
\text{B: } & 10^{-3} \times 0.01 \times 1.0 = 10^{-5}
\end{align*}
Under L2-in-Adam, B shrinks a thousand times faster than A, for no reason connected to the
regularisation you asked for. Under AdamW both shrink equally, as the coefficient $\lambda_w$
intends.
\end{worked}

\section{Learning-rate schedules}\label{sec:trainloop:schedule}

A fixed learning rate is a compromise: large steps cover ground fast but cannot settle into a
narrow minimum; small steps settle but crawl. A \emph{schedule}\index{learning-rate schedule}
changes $\eta$ over the run. Two shapes matter here.

\subsection{Linear warmup}\label{sec:trainloop:warmup}

\begin{definition}[Linear warmup]\index{warmup}
Over the first $W$ epochs (or steps), raise the learning rate linearly from a small $\eta_{\min}$
to the target $\eta_{\max}$:
\[
\eta(t) = \eta_{\min} + (\eta_{\max} - \eta_{\min})\,\frac{t}{W}, \qquad 0 \le t \le W.
\]
\end{definition}
With $W = 5$, $\eta_{\min} = 10^{-6}$ and $\eta_{\max} = 10^{-3}$ the five warmup epochs would use
$10^{-6}$, $2.008 \times 10^{-4}$, $4.006 \times 10^{-4}$, $6.004 \times 10^{-4}$ and $8.002
\times 10^{-4}$, reaching $10^{-3}$ at $t = 5$. (Row $t = 1$: $10^{-6} + 0.000999 \times 1/5 =
0.000001 + 0.0001998 = 0.0002008$.)

Why warm up? \Cref{we:trainloop:adam} showed that Adam's first steps move every parameter by
about $\eta$, steered by moment estimates built from one or two batches, on a network whose
parameters are random. Large, poorly aimed early steps can push the network into a bad region
it never recovers from. Warmup keeps those steps small until the moment estimates have seen
enough batches to be trustworthy. \citet{goyal2017accurate} made warmup standard for large-batch
training.

\begin{sourcenote}[: \texttt{warmup\_epochs} is documented, defaulted, set in every config, and never read]
\file{app/models/training/training_config.py}, lines 418--422, defines
\begin{lstlisting}
warmup_epochs: int = Field(default=5, ge=0,
    description="Linear warmup from min_lr to learning_rate")
\end{lstlisting}
All 22 training configurations in \file{configs/} set it (both Indradhanu configurations to 5),
the fine-tuning configurations set it to 0 or 1 and the SpectralEarth pretraining configurations to
0. A search of every Python file in the repository,
\begin{lstlisting}[language=bash]
grep -rn -i "warmup" --include="*.py" .
\end{lstlisting}
returns three lines, all in \file{training_config.py}: the field (418), its description (421)
and a docstring example (505). \code{_build_scheduler}
(\file{app/abstract_classes/foundation_trainer.py}, lines 517--553) builds \code{cosine},
\code{step} or \code{plateau} and consults none of it. Checked 2026-10-29; the teaching skill's
note is still true. Every run in this repository, including the v0.2.0 run, started at the full
learning rate on step 1.

The course (IW28) links this to the v0.2.0 history: ``v0.2.0 epoch 2 hits 0.035 (lr=1e-3 was too
aggressive; later epochs regress to \textasciitilde 0.04)''
\tagM{research/model_break_down/05_hyperspectral_segformer_mae.md}. An unusually good early
epoch followed by a regression that never fully recovers is consistent with early
instability, the thing warmup exists to prevent; it is not proof. The lesson the repository
states as policy: verify a config field against the source before assuming it does anything.
\end{sourcenote}

\subsection{Cosine annealing}

\begin{definition}[Cosine annealing]\index{cosine annealing}
Over $T$ epochs, from $\eta_{\max}$ down to $\eta_{\min}$:
\begin{equation}\label{eq:trainloop:cosine}
\eta(t) = \eta_{\min} + \frac{1}{2}\,(\eta_{\max} - \eta_{\min})\left(1 + \cos\frac{\pi t}{T}\right), \qquad 0 \le t \le T.
\end{equation}
\end{definition}
This is the schedule of \citet{loshchilov2017sgdr} without restarts, and the PyTorch class
\code{CosineAnnealingLR} implements it. Check the endpoints: at $t = 0$, $\cos 0 = 1$ and the
bracket is 2, so $\eta = \eta_{\min} + (\eta_{\max} - \eta_{\min}) = \eta_{\max}$. At $t = T$,
$\cos\pi = -1$, the bracket is 0, and $\eta = \eta_{\min}$.

\paragraph{The shape.} Differentiate with respect to $t$ (chain rule: the derivative of
$\cos(\pi t/T)$ is $-\sin(\pi t/T)\cdot\pi/T$):
\[
\frac{d\eta}{dt} = -\frac{\pi}{2T}\,(\eta_{\max} - \eta_{\min})\,\sin\frac{\pi t}{T}.
\]
The slope is zero at $t = 0$ and $t = T$ (where $\sin$ is 0) and steepest at $t = T/2$. So the
learning rate stays high and nearly flat at the start, falls fastest in the middle, and
flattens out near the floor at the end.

\begin{worked}[label={we:trainloop:cosine}]{The v0.2.0 schedule at five epochs}
v0.2.0 uses $\eta_{\max} = 10^{-3}$, $\eta_{\min} = 10^{-6}$ and $T = 200$
\tagM{configs/hyperspectral_segformer_exp_2.json}. So $\frac12(\eta_{\max} - \eta_{\min}) =
\frac12 \times 0.000999 = 0.0004995$.
\begin{center}
\begin{tabular}{rrrrl}
\toprule
$t$ & $\pi t/200$ & $\cos$ & $1 + \cos$ & $\eta(t)$\\
\midrule
0 & 0 & 1.000000 & 2.000000 & $10^{-6} + 0.0004995 \times 2 = 1.000 \times 10^{-3}$\\
50 & $\pi/4$ & 0.707107 & 1.707107 & $10^{-6} + 0.0004995 \times 1.707107 = 8.537 \times 10^{-4}$\\
100 & $\pi/2$ & 0.000000 & 1.000000 & $10^{-6} + 0.0004995 \times 1 = 5.005 \times 10^{-4}$\\
150 & $3\pi/4$ & $-0.707107$ & 0.292893 & $10^{-6} + 0.0004995 \times 0.292893 = 1.473 \times 10^{-4}$\\
200 & $\pi$ & $-1.000000$ & 0.000000 & $10^{-6} + 0 = 1.000 \times 10^{-6}$\\
\bottomrule
\end{tabular}
\end{center}
Row $t = 100$ in full: $0.0004995 \times 1 = 0.0004995$; $0.000001 + 0.0004995 = 0.0005005$.
PyTorch's \code{CosineAnnealingLR}, stepped as the trainer steps it, gives the same five values
to the digits shown \tagC.
\end{worked}

\paragraph{When the schedule steps.} The trainer calls \code{self.scheduler.step()} once, at the
end of each epoch, after validation (\file{app/abstract_classes/foundation_trainer.py}, lines
237--243). So epoch index $e$ trains at $\eta(e)$: index 0 at $10^{-3}$, the last index, 199, at
$\eta(199) = 1.06 \times 10^{-6}$ \tagC, and $\eta(200) = 10^{-6}$ is reached only after training
ends. The log line is written \emph{after} the step, so the ``lr'' it prints for logged epoch
$k$ (index $k - 1$) is $\eta(k)$, the rate the \emph{next} epoch will use.

\begin{sourcenote}[: the example log line's learning rate]
The course (IW28) shows a log line ``Epoch 101/200 \ldots{} lr: 4.94e-04''. Logged epoch 101
prints $\eta(101)$: $\cos(101\pi/200) = -0.015707$, so $\eta = 10^{-6} + 0.0004995 \times
0.984293 = 10^{-6} + 0.00049165 = 4.93 \times 10^{-4}$ (PyTorch: $4.9265 \times
10^{-4}$ \tagC). The course's 4.94e-4 is an illustrative line, off in the third digit.
\end{sourcenote}

\begin{figure}[htbp]
\centering
\begin{tikzpicture}
\begin{axis}[width=0.92\linewidth, height=6cm, xmin=0, xmax=200, ymin=0, ymax=1.05,
  xlabel={epoch index $t$}, ylabel={learning rate ($\times 10^{-3}$)}, grid=major,
  grid style={gray!20}, xtick={0,25,50,75,100,125,150,175,200}, legend pos=north east,
  legend style={font=\small}]
\addplot[loss, very thick, domain=0:200, samples=201] {0.001 + 0.4995*(1 + cos(180*x/200))};
\addlegendentry{v0.2.0: $10^{-3} \rightarrow 10^{-6}$, $T = 200$}
\addplot[spat, very thick, dashed, domain=0:200, samples=201] {0.001 + 0.2495*(1 + cos(180*x/200))};
\addlegendentry{SpectralEarth F001: $5\times 10^{-4} \rightarrow 10^{-6}$, $T = 200$}
\addplot[only marks, mark=*, mark size=1.8pt, loss, forget plot] coordinates
  {(0,1.0) (50,0.8537) (100,0.5005) (150,0.1473) (200,0.001)};
\addplot[maskgrey, thick, domain=0:5, samples=2] {0.001 + (0.999)*x/5};
\addlegendentry{a 5-epoch linear warmup (never implemented)}
\end{axis}
\end{tikzpicture}
\caption[The learning-rate schedules of v0.2.0 and F001]{The learning-rate schedules of two real configurations, drawn to scale from
\cref{eq:trainloop:cosine}: v0.2.0 (\texttt{configs/hyperspectral\_segformer\_exp\_2.json}) and the
SpectralEarth final pretraining F001
(\texttt{research/experiments/indradhanu-se-pretrain/F001-final-pretrain/config.json}). Dots are
the worked values of \cref{we:trainloop:cosine}. The grey ramp at the left is what
\texttt{warmup\_epochs = 5} describes and what no code implements.}
\label{fig:trainloop:cosine}
\end{figure}

\begin{drift}
Learning-rate settings as of 2026-10-29: v0.2.0, $10^{-3}$ over 200 epochs; v0.1.0, $10^{-4}$
over 500; SpectralEarth pretraining F001, $5 \times 10^{-4}$ over 200
\tagM{research/experiments/indradhanu-se-pretrain/F001-final-pretrain/config.json}; fine-tuning
runs between $3 \times 10^{-5}$ and $10^{-3}$ over 3 to 20 epochs (FTH3, the run shipped as
Indradhanu V2 v2.1.0, used $3 \times 10^{-5}$ over 6 epochs
\tagM{research/experiments/indradhanu-finetune/FTH3-heavy-joint-more2/config.json}). All of them
use \code{min_lr} $= 10^{-6}$ and the cosine schedule.
\end{drift}

\paragraph{A resume bug, told as a lesson.} When a run is resumed, the scheduler must continue
from where it stopped. The original resume code rebuilt the scheduler with
\code{last_epoch=k}, but \code{CosineAnnealingLR} computes each rate recursively from the
optimiser's current rate, so the rebuild took one extra step: a run killed after epoch 2
continued at $10^{-6}$ and then $1.47 \times 10^{-4}$ instead of $1.47 \times 10^{-4}$ and then
$10^{-6}$ (measured 2026-10-14, \file{app/abstract_classes/foundation_trainer.py}, lines
580--584). The fix saves and restores the scheduler's own state. A schedule is part of the
model's history, and resuming it is as delicate as resuming the weights.

\section{Gradient clipping}\label{sec:trainloop:clip}

\begin{definition}[Gradient clipping by norm]\index{gradient clipping}
If the gradient's length (L2 norm over all parameters) exceeds a threshold $c$, scale the whole
gradient down to length $c$; otherwise leave it:
\[
\vect{g} \leftarrow \vect{g}\cdot\min\!\left(1, \frac{c}{\norm{\vect{g}}}\right).
\]
\end{definition}
Clipping by norm keeps the gradient's direction and bounds only the step's size, a guard
against a single bad batch (a glitch, a NaN-adjacent value) throwing the parameters far away
\citep{pascanu2013difficulty}.

\begin{worked}[label={we:trainloop:clip}]{Clipping a 3-vector}
$\vect{g} = (3, 4, 12)$, threshold $c = 1$.
\begin{align*}
\norm{\vect{g}}^2 &= 3^2 + 4^2 + 12^2 = 9 + 16 + 144 = 169\\
\norm{\vect{g}} &= \sqrt{169} = 13\\
\min(1, 1/13) &= 1/13 = 0.076923\\
\vect{g}_{\text{clipped}} &= (3/13,\ 4/13,\ 12/13) = (0.230769,\ 0.307692,\ 0.923077)
\end{align*}
Check the length: $0.230769^2 + 0.307692^2 + 0.923077^2 = 0.053254 + 0.094675 + 0.852071 =
1.000000$. The direction is unchanged: the ratios $3 : 4 : 12$ survive. PyTorch's
\code{clip_grad_norm_} returns $(0.230769, 0.307692, 0.923077)$ \tagC{} (it divides by
$\norm{\vect{g}} + 10^{-6}$, which changes the seventh digit).
\end{worked}

\begin{sourcenote}[: no clipping]
A grep for \code{clip_grad} across \file{app/}, \file{scripts/} and \file{research/} finds
nothing: no trainer in this repository clips gradients. The loss's own guards
(\cref{sec:losses:atan2}) and the 40\% patch filter are the only protections against a wild
batch. The course (IW28) says ``no gradient clipping''; correct.
\end{sourcenote}

\section{Mixed precision}\label{sec:trainloop:precision}

A floating-point number stores a sign, an exponent (the range) and a mantissa (the
precision). Three formats matter:
\begin{center}
\begin{tabular}{lrrll}
\toprule
format & exponent bits & mantissa bits & relative spacing & range\\
\midrule
float32 & 8 & 23 & $2^{-23} \approx 1.2 \times 10^{-7}$ & $\sim 10^{\pm 38}$\\
bfloat16 & 8 & 7 & $2^{-7} \approx 0.0078$ & $\sim 10^{\pm 38}$\\
float16 & 5 & 10 & $2^{-10} \approx 0.00098$ & $\sim 6.5\times 10^{4}$ max\\
\bottomrule
\end{tabular}
\end{center}
\emph{Mixed precision}\index{mixed precision}\index{bfloat16} runs the expensive matrix
multiplications in a 16-bit format (half the memory, much faster on modern GPUs) and keeps
precision-critical parts in float32 \citep{micikevicius2018mixed}. bfloat16 keeps float32's
exponent, so it has the same range and never underflows where float32 would not; float16 needs
\emph{loss scaling} to keep small gradients from flushing to zero, bfloat16 does not.

\begin{worked}[label={we:trainloop:bf16}]{Why the loss stays in float32}
A reflectance near 0.2 lies in $[0.125, 0.25) = [2^{-3}, 2^{-2})$, where bfloat16's spacing is
$2^{-3} \times 2^{-7} = 2^{-10} = 0.000977$. So
\begin{align*}
0.2000 &\to 0.2001953 \text{ in bfloat16}\\
0.2003 &\to 0.2001953\\
0.2005 &\to 0.2001953
\end{align*}
(PyTorch's conversions \tagC). A reconstruction error of $0.0003$ or $0.0005$, the size a good
model's L1 is made of (\cref{ch:losses}), vanishes completely. In float32 the three values stay
distinct.
\end{worked}

The SpectralEarth and fine-tuning trainers wrap only the model's forward pass in bfloat16
autocast and convert its output back before the loss:
\begin{lstlisting}
bf16 = self.config.model_config_.amp_bf16 and self.device.type == "cuda"
with torch.autocast("cuda", dtype=torch.bfloat16, enabled=bf16):
    out = model(pixels, wl, band_valid, keep_mask, out_wl)
return out.float()          # the loss - small differences of reflectance - in full precision
\end{lstlisting}
(\file{app/foundation_models/trainers/spectral_stem_mae_trainer.py}, lines 269--275). Every
SpectralEarth configuration except the laptop smoke tests, and every fine-tuning configuration,
sets \code{amp_bf16: true}.\newline
\tagM{research/experiments/indradhanu-se-pretrain/F001-final-pretrain/config.json}\newline
The configuration field's description says it roughly halves the cross-attention stem's activation
memory. Be precise about what the cast buys. The prediction is still computed by bfloat16 arithmetic
inside the network, so it carries bfloat16-sized rounding in the network's own (normalised)
units; what \code{out.float()} guarantees is that the truth $\vect{x}$ is never rounded and that
the subtraction, the absolute values, the angle and the averaging over millions of entries are
done in float32. There is no gradient scaler (\code{GradScaler}) anywhere, which is correct for
bfloat16. The original v0.2.0 trainer (\file{app/foundation_models/trainers/hyperspectral_segformer_mae_trainer.py})
has no autocast: it trained entirely in float32.

\section{Gradient accumulation}\label{sec:trainloop:accum}

When the batch you want does not fit in GPU memory, process $K$ smaller batches and step once:
\begin{lstlisting}
scaled_loss = loss / accum_steps          # divide so the sum of K gradients is their mean
scaled_loss.backward()                    # .grad += d(scaled_loss)/dw
if accum_count % accum_steps == 0:
    self.optimizer.step()
    self.optimizer.zero_grad()
\end{lstlisting}
(\file{app/foundation_models/trainers/hyperspectral_segformer_mae_trainer.py}, lines 100--148.)
Because \code{backward()} adds into \code{.grad}, after $K$ calls the field holds
$\sum_{k=1}^{K}\nabla(L_k/K) = \frac{1}{K}\sum_k \nabla L_k$, the mean of the $K$ batch
gradients. If each batch has $B$ patches and $L_k$ is the mean over them, this is
\[
\frac{1}{K}\sum_{k=1}^{K}\frac{1}{B}\sum_{i \in \text{batch }k}\nabla\ell_i = \frac{1}{KB}\sum_{\text{all }KB}\nabla\ell_i,
\]
exactly the gradient of one batch of $KB$ patches. Without the division by $K$ the gradient
would be $K$ times larger; with Adam that barely matters (it divides by the gradient's size), but
with plain descent it would multiply the learning rate by $K$. If the loader runs out mid-group,
the leftover gradients are applied (``the flush'', lines 143--146), not discarded.

Two cautions. The equality needs equal batch sizes; after the 40\% filter, batches carry
different numbers of surviving patches, so the accumulated gradient weights each batch equally
rather than each patch. And accumulation does not reproduce a large batch's BatchNorm
statistics, which are still computed per small batch (\cref{ch:finetune}).

\begin{drift}
v0.2.0 set \code{gradient_accumulation_steps: 1}, so the override returns straight to the base
loop. The SpectralEarth pretraining F001 uses \code{batch_size} 8 and accumulation 4, an
effective batch of $8 \times 4 = 32$ and at most $30{,}000/32 = 937.5$, so 938, optimiser steps
per epoch \tagM{research/experiments/indradhanu-se-pretrain/F001-final-pretrain/config.json}.
\end{drift}

\section{The training data}\label{sec:trainloop:data}

\coursefig[0.95\linewidth]{iw-27-data}{The training-data pipeline, one job per stage. A raw scene
becomes a 165-band vendable; PatchPlanGenerator decides where to cut 128 $\times$ 128 patches
(stride 64 in training, so neighbours overlap by half, the last snapping to the scene edge);
the patcher cuts them; they are packed into .tar shards on S3 and streamed by the DataLoader as
batches.}{fig:trainloop:data}

\subsection{Deciding where to cut}

\code{PatchPlanGenerator} (\file{app/utils/patch_generation/generate_patch_plan.py}) returns
coordinates only; no pixels are touched. Given the cube's height and width, a patch size and a
\emph{stride}\index{stride} (how far to move between patches), it lists every top-left corner.
The training convention is stride $=$ size$/2$, so 128-pixel patches step by 64 and neighbours
overlap by half: every interior pixel appears in several patches, sometimes centred, sometimes
near an edge, which evens out the poorer reconstruction near patch borders.

When the next step would run off the edge, the last patch is snapped back to end flush with
the boundary (lines 66--79):
\begin{lstlisting}
if row + request.height > cube_height:
    row_coords.append(cube_height - request.height)   # last patch ends exactly at the edge
    break
\end{lstlisting}

\begin{worked}[label={we:trainloop:plan}]{The row coordinates of a 300-pixel scene}
Height 300, size 128, stride 64. Walk the loop.
\begin{align*}
\text{row } 0:&\quad 0 + 128 = 128 \le 300, \text{ keep } 0\\
\text{row } 64:&\quad 64 + 128 = 192 \le 300, \text{ keep } 64\\
\text{row } 128:&\quad 128 + 128 = 256 \le 300, \text{ keep } 128\\
\text{row } 192:&\quad 192 + 128 = 320 > 300, \text{ snap to } 300 - 128 = 172, \text{ stop}
\end{align*}
Rows $[0, 64, 128, 172]$: four, and the last gap is $172 - 128 = 44$, not 64. The patch count is
not $\lceil 300/64 \rceil = 5$; the module's docstring warns about exactly this. The real
generator returns $[0, 64, 128, 172]$ \tagC.
\end{worked}

\begin{sourcenote}[: a duplicated last patch]
Run the same loop on a height of 256: rows 0, 64, 128 are kept ($128 + 128 = 256 \le 256$); at
row 192, $192 + 128 = 320 > 256$, so it snaps to $256 - 128 = 128$, a coordinate it already
has. The real generator returns rows $[0, 64, 128, 128]$, and for a dimension of exactly 128 it
returns $[0, 0]$ \tagC{} (\file{research/textbook/tools/compute_24-trainloop_output.txt}). The
duplicate appears whenever $(\text{extent} - \text{size})$ is a multiple of the stride, because
the loop continues while \code{row < cube_height} rather than stopping once a patch has
reached the edge. The patcher (\file{app/utils/patch_generation/hyperspectral_patcher.py},
line 96) iterates the plan without removing duplicates, so such a scene emits the same patch
twice under the same key. That is consistent with the ``duplicate file name in tar file''
webdataset warning that the hyperspectral trainer suppresses
(\file{app/foundation_models/trainers/hyperspectral_segformer_mae_trainer.py}, line 21, whose
comment attributes it to resampled shards); whether the shards actually contain such
duplicates was not checked here. The effect on training is a slight over-weighting of edge
patches; at inference, an overlapping patch predicted twice is harmless if overlaps are
averaged.
\end{sourcenote}

\subsection{Doing the cutting, and the shard contract}

For each coordinate the patcher emits one dictionary, written as named members of a tar file:
\begin{lstlisting}
{"__key__":           f"{scene_id}#row_coord:{r}#col_coord:{c}",
 "meta.json":         {... "patch_height": 128, "patch_stride": 64, "band_count": 165 ...},
 "pixels.npy":        cube[:, r:r+h, c:c+w].astype(np.float32),      # (165, 128, 128)
 "validity_cube.npy": validity[:, r:r+h, c:c+w].astype(np.int8),     # (165, 128, 128) 0/1
 "wavelengths.npy":   wavelengths}                                   # (165,)
\end{lstlisting}
The key names are a contract: the trainer reads \code{batch["pixels.npy"]} and
\code{batch["validity_cube.npy"]} literally. The mask is int8 because it only ever holds 0 or 1,
a quarter of float32's size.

\begin{definition}[Shard]\index{shard}
A \emph{shard} is a \texttt{.tar} archive packing many patches, read sequentially. Millions of
tiny files would each cost a network request; a few hundred large shards stream efficiently.
This is the \emph{webdataset} convention, which is why the keys look like file names.
\end{definition}

The S3 location is a template,
\begin{lstlisting}[language={}]
patches/{provider}/{split}/{stage}/w{size}_h{size}_s{stride}/
\end{lstlisting}
filled for Indradhanu as \code{patches/hyperspectral/train/final/w128_h128_s64/} in bucket
\code{allotrope-raw-data-india} \tagM{configs/hyperspectral_segformer_exp_2.json}. The provider is
``hyperspectral'', not ``prisma'' or ``enmap'': PRISMA and EnMAP patches share the same
shards, interleaved, because on the 165-band common grid (\cref{ch:sensors}) they are
indistinguishable at the tensor level. The schema's default template points at Landsat thermal
data, and the source says so (\file{app/models/training/training_config.py}, line 296):
\begin{lstlisting}
# Important to set this properly The default is landsat which is dangerous.
\end{lstlisting}

\paragraph{Hot storage.} Shards are copied once to the training machine's local disk and kept
(\code{hot_storage}: 500 training and 100 test shards for v0.2.0, \code{skip_sync_if_exists:
true}). There is no rotation in the base trainer: those 500 shards are the training set for the
whole run, and variety across epochs comes from \code{shardshuffle} (30 in v0.2.0), which
reorders shards and samples within a buffer.

\subsection{What an epoch means here}\label{sec:trainloop:epoch}

The corpus is too large to pass over even once in reasonable time, so an epoch is a sample
budget:
\begin{lstlisting}[language={}]
"train_samples_per_epoch": { "128": 40000 },
"test_samples_per_epoch":  { "128":  5000 }
\end{lstlisting}
``Consume 40,000 training patches of size 128, then 5,000 validation patches, and call that an
epoch.'' The dictionary is keyed by patch size because the thermal models train on several sizes
at once; Indradhanu uses only 128. Only patches that survive the 40\% filter
(\cref{sec:losses:forty}) count:
\begin{lstlisting}
for batch in loader:
    if valid_samples >= sample_cap:      # 40,000
        break
    loss, num_kept = self.compute_loss(batch, self.model)
    if num_kept == 0:
        continue                         # whole batch filtered out: no step, no budget used
    ...
    valid_samples += num_kept
\end{lstlisting}
With $B = 128$ and no filtering, $40{,}000/128 = 312.5$, so the loop runs 313 batches and stops
at $313 \times 128 = 40{,}064$ samples \tagC: the cap is checked before each batch, so the last
batch overshoots by up to $B - 1$. Over 200 epochs that is about $40{,}000 \times 200 =
8{,}000{,}000$ patch presentations \tagC, drawn from the same 500 cached shards.

\paragraph{Feeding it is an I/O problem.} One patch is $165 \times 128 \times 128$ float32 values,
\begin{align*}
165 \times 128 &= 21{,}120\\
21{,}120 \times 128 &= 2{,}703{,}360 \text{ values}\\
2{,}703{,}360 \times 4 \text{ bytes} &= 10{,}813{,}440 \text{ bytes} \approx 10.8 \text{ MB},
\end{align*}
and a batch of 128 is $1{,}384{,}120{,}320$ bytes, about 1.4 GB, before any activation \tagC.
That is why v0.2.0 ran 40 loader workers (\code{num_workers}). If the data pipeline stalls, an
expensive GPU sits idle.

\subsection{Randomness and augmentation}

The v0.2.0 trainer applies no geometric augmentation: no flips, rotations or crops (a search of
the trainers for \code{flip} and \code{rot90} finds none). The variety each epoch comes from
three random sources: the shard and sample order (\code{shardshuffle}); the token mask, redrawn
for every batch so that a different 65\% of each patch's valid tokens is hidden each time it
is seen (\cref{ch:masking}); and, in validation, the random half-split of
\cref{ch:losses}.

The later trainers add two augmentations, each taught where it belongs:
\begin{itemize}
\item \textbf{Band-list augmentation} (SpectralEarth pretraining and fine-tuning;
\code{sample_band_subset} in \file{app/foundation_models/trainers/spectral_stem_mae_trainer.py},
lines 75--92). With probability \code{band_augment_prob} (0.8 in F001) a batch's input is a
random subset of its bands, in one of three shapes: thinning (keep 50--95\%), one to three gaps of
20--150 nm, or the bands nearest an even grid of 6--15 nm spacing, never fewer than 40 bands; the
stated wavelengths are jittered by up to $\pm$\code{band_jitter_nm} (3 nm in F001), and the loss
still covers every band. It teaches a wavelength-aware stem to cope with sensors it was not
trained on (\cref{ch:crossstem,ch:pretrain}).
\item \textbf{Implant augmentation} (fine-tuning run FT05): synthetic objects implanted in the
input, the clean scene as the target (\cref{ch:ftruns}).
\end{itemize}
The SpectralEarth shards also change the data contract: \code{pixels.npy} is int16, reflectance
$\times 10{,}000$, with $-32768$ marking a dead voxel, and validity is per band rather than per
pixel (\cref{ch:pretrain}).

\section{The loop}\label{sec:trainloop:loop}

\subsection{From the command line to the epoch loop}

\begin{lstlisting}[language=bash]
python scripts/train_foundation_model.py configs/hyperspectral_segformer_exp_2.json
\end{lstlisting}
The script validates the JSON into a pydantic \code{TrainingConfig} (a typo fails at once, not
forty minutes in), looks up the trainer class by model name in
\file{app/foundation_models/trainers/trainer_factory.py}, and calls \code{trainer.train()}.

\code{FoundationTrainer} (\file{app/abstract_classes/foundation_trainer.py}) owns everything
generic: device, optimiser, scheduler, data loaders, hot storage, the epoch loop, checkpoints and
logging. A model's trainer implements three methods, \code{build_model}, \code{compute_loss}
(returns the loss and the number of surviving samples) and \code{validation_step}; Indradhanu's
also overrides \code{_run_epoch} (to record the epoch index for the SAM ramp and log the loss
components) and \code{_run_train_pass} (to add gradient accumulation).

\begin{figure}[htbp]
\centering
\begin{tikzpicture}[node distance=5mm and 8mm, font=\small,
  box/.style={draw, rounded corners=2pt, align=center, minimum height=8mm, text width=37mm, fill=white},
  dec/.style={draw, diamond, aspect=2.4, align=center, inner sep=1pt, fill=boxwarn, font=\scriptsize},
  arr/.style={-{Stealth[length=2mm]}, thick}]
\node[box, fill=gray!8] (cfg) {config $\rightarrow$ trainer\\ build model, \textbf{Adam}, \textbf{cosine}};
\node[box, below=of cfg] (train) {\texttt{model.train()}};
\node[dec, below=of train] (cap) {surviving samples\\ $\ge$ 40,000?};
\node[box, below=of cap, fill=spec!10] (batch) {next batch \shape{(B,165,128,128)}\\ drop patches $<$ 40\% valid};
\node[box, below=of batch, fill=spat!8] (fwd) {random token mask;\\ forward pass $\rightarrow \xhat$};
\node[box, below=of fwd, fill=loss!12] (loss) {$L = L_1 + \lambda(t)\SAM$\\ on hidden, valid, eroded pixels};
\node[box, below=of loss, fill=attn!10] (bwd) {\texttt{zero\_grad}; \texttt{backward};\\ \texttt{optimizer.step()}};
\node[box, right=22mm of cap] (val) {\texttt{model.eval()}, no grad:\\ two-pass validation\\ to 5,000 samples};
\node[box, below=of val] (sched) {\texttt{scheduler.step()}\\ $\eta(t) \rightarrow \eta(t+1)$};
\node[box, below=of sched] (log) {log $L$, $L_1$, $\SAM$, $\lambda$, lr};
\node[box, below=of log] (ckpt) {checkpoint every $n$ epochs;\\ keep top $k$ by val loss};
\node[dec, below=of ckpt] (more) {more epochs?};
\node[box, below=of more, fill=gray!8] (done) {done};
\draw[arr] (cfg) -- (train);
\draw[arr] (train) -- (cap);
\draw[arr] (cap) -- node[right, font=\scriptsize]{no} (batch);
\draw[arr] (batch) -- (fwd);
\draw[arr] (fwd) -- (loss);
\draw[arr] (loss) -- (bwd);
\draw[arr] (bwd.west) -- ++(-6mm,0) |- node[pos=0.25, left, font=\scriptsize, align=right]{next\\ step} (cap.west);
\draw[arr] (cap) -- node[above, font=\scriptsize]{yes} (val);
\draw[arr] (val) -- (sched);
\draw[arr] (sched) -- (log);
\draw[arr] (log) -- (ckpt);
\draw[arr] (ckpt) -- (more);
\draw[arr] (more) -- node[right, font=\scriptsize]{no} (done);
\coordinate (rx) at ($(val.east)+(6mm,0)$);
\draw[arr] (more.east) -- (more.east -| rx) -- node[pos=0.5, right, font=\scriptsize, align=left]{next epoch,\\ $t \leftarrow t+1$} (train.east -| rx) -- (train.east);
\end{tikzpicture}
\caption{The training loop of \texttt{FoundationTrainer} with Indradhanu's loss, as the source runs
it. The left column is one training step, repeated until the epoch's sample budget is spent;
the right column runs once per epoch: validation, the scheduler step, logging and the periodic
checkpoint. The SAM weight $\lambda(t)$ and the learning rate $\eta(t)$ both change only at the
epoch boundary.}
\label{fig:trainloop:flow}
\end{figure}

\coursefig[0.9\linewidth]{iw-28-loop}{The course's drawing of one epoch: a training pass with
random masking and gradient accumulation, a validation pass, a scheduler step, logging, and a
periodic checkpoint. The dashed frame repeats once per epoch.}{fig:trainloop:courseloop}

\subsection{What one epoch does}

\begin{lstlisting}
def _run_epoch(self, epoch, train_loaders, test_loaders):
    self.model.train()
    for size in patch_sizes:                                  # Indradhanu: only 128
        size_loss, size_samples = self._run_train_pass(train_loaders[size], cap)
    avg_train_loss = epoch_train_loss / max(epoch_train_samples, 1)
    self.model.eval()
    for size in patch_sizes:
        val_losses[size] = self._run_val_pass(test_loaders[size], cap)   # under torch.no_grad()
    avg_val_loss = sum(val_losses.values()) / len(val_losses)
    self.scheduler.step()
    ...log...
    if (epoch + 1) % ckpt.save_every_n_epochs == 0:
        self._save_checkpoint(epoch + 1, avg_train_loss, val_losses)
        self._cleanup_checkpoints()
\end{lstlisting}
Five phases: train, validate, step the scheduler, log, maybe checkpoint
(\file{app/abstract_classes/foundation_trainer.py}, lines 202--274).

\paragraph{\texttt{model.train()} and \texttt{model.eval()}} do not train or evaluate anything.
They flip a flag that changes two kinds of layer. Dropout zeroes random activations in train
mode and does nothing in eval mode; BatchNorm uses the current batch's statistics and updates
its running averages in train mode, and uses the stored averages in eval mode. Indradhanu has
both (\code{drop_rate} 0.4 and a BatchNorm in the compressor), so forgetting \code{eval()} before
validation would make validation noisy and slightly optimistic without any error. Do not
confuse \code{model.train()} (a mode flag on the network) with \code{trainer.train()} (the
method that runs the whole loop).

\paragraph{The training step} (lines 276--307) is the move of the opening box:
\code{optimizer.zero_grad()}, \code{compute_loss}, \code{loss.backward()},
\code{optimizer.step()}.

\subsection{Sample-weighted averaging}

\code{compute_loss} returns a per-sample mean, so the pass multiplies it back by the number of
surviving samples before summing:
\begin{lstlisting}
total_loss += loss.item() * num_kept
valid_samples += num_kept
...
return total_loss / max(valid_samples, 1)        # max(.., 1): an all-filtered epoch gives 0, not a crash
\end{lstlisting}

\begin{worked}[label={we:trainloop:weighted}]{Why the averaging is weighted}
Two batches: in the first only 3 patches survived the filter, with mean loss 0.05; in the
second all 128 survived, with mean loss 0.03.

\textbf{Weighted} (what the code does):
\begin{align*}
0.05 \times 3 &= 0.15\\
0.03 \times 128 &= 3.84\\
(0.15 + 3.84)/(3 + 128) &= 3.99/131 = 0.030458
\end{align*}
\textbf{Unweighted} (averaging the two batch means): $(0.05 + 0.03)/2 = 0.04$.

The unweighted answer lets 3 patches count as much as 128. The code carries a comment
explaining the multiplication, because it is the kind of line a reader might ``simplify'' away.
\end{worked}

\subsection{Checkpoints and resuming}

Each checkpoint stores the model's parameters and buffers, the optimiser's state (Adam's $m$ and
$v$ for every parameter, needed to resume cleanly), the scheduler's state, the epoch, the
training and per-size validation losses, and the whole configuration. Files are named
\texttt{\{model\_name\}\_v\{version\}\_epoch\{epoch\}.pt}, which is why the shipped v0.2.0 file is
\code{hyperspectral_segformer_mae_v0.2.0_epoch200.pt}; \code{keep_top_k} (20 in v0.2.0) prunes
all but the best by validation loss.

\begin{center}
\begin{tabular}{p{2.2cm}p{5.6cm}p{4.6cm}}
\toprule
\code{resume_mode} & restores & use it when\\
\midrule
\code{"resume"} & weights, optimiser state, epoch counter, scheduler state & a run crashed or was stopped\\
\code{"finetune"} & weights only; fresh optimiser, scheduler, epoch 0 & adapting a trained model to new data\\
\bottomrule
\end{tabular}
\end{center}
Resuming a 200-epoch cosine run at epoch 150 with a fresh scheduler would restart the learning
rate at $10^{-3}$ and undo the decay; \code{"resume"} exists to prevent that. The SpectralEarth
trainer adds \code{last.pt} and \code{best.pt}, mirrored to S3, automatic resumption, and early
stopping on the validation score (\cref{ch:pretrain}).

\begin{sourcenote}[: which number ranks the checkpoints]
\code{keep_top_k} and \code{best.pt} rank epochs by validation loss. In the base Indradhanu trainer
that loss uses the ramped SAM weight (\cref{ch:losses}), so epochs inside the ramp are on a
different scale from later ones. The SpectralEarth trainer found this the hard way (an
early-stop rule that would have fired whatever the model did) and now scores validation at the
full weight from the first epoch; \cref{ch:pretrain} tells that story.
\end{sourcenote}

\subsection{Reading a log line}

\begin{lstlisting}[language={}]
Epoch 101/200 | train_loss: 0.041203 | val_loss: [128px: 0.043490] | avg_val: 0.043490 | lr: 4.93e-04
  L1: 0.009812 | SAM: 0.062776 rad (3.60 deg) | lambda: 0.500
\end{lstlisting}
(An illustrative line in the format of lines 245--256 of the base trainer and 65--68 of
Indradhanu's, with the learning rate corrected as in the source check above.)
\code{train_loss} and \code{val_loss} are the combined $L_1 + \lambda\SAM$; \code{lr} is the rate
the next epoch will use; the second line gives the components. Watch the second line, for the
reasons in \cref{ch:losses}.

\begin{keypoint}
The loop is one move repeated: batch, loss, gradient, Adam step; the epoch boundary is where the
learning rate, the SAM weight, validation and checkpoints change, and an epoch here is a
sample budget, not a pass over the data.
\end{keypoint}

\begin{confusions}
\confusion{Does model.train() do the training?}{No. It sets a mode flag on the network (dropout
on, BatchNorm in batch-statistics mode). \code{trainer.train()}, a different method on a
different object, runs the loop.}
\confusion{Is an epoch one pass over the training data?}{Not here. It is a budget of 40,000
surviving training patches and 5,000 validation patches (v0.2.0), drawn from 500 cached shards.
Coverage is statistical, not exhaustive.}
\confusion{Does warmup\_epochs = 5 mean the run warmed up?}{No. The field is never read; every
run started at the full learning rate (\cref{sec:trainloop:warmup}).}
\confusion{Is Adam with weight decay the same as AdamW?}{No. Adding an L2 term to the gradient
sends the decay through Adam's normaliser, so weights with large gradients decay less; AdamW
subtracts the decay separately (\cref{we:trainloop:adamw}). This project uses neither: weight
decay is 0.}
\confusion{Why is the learning rate in the log not the one the epoch trained at?}{The scheduler
steps before the log line is written, so the printed rate is the next epoch's.}
\confusion{Can I change the learning rate mid-run?}{Edit the config and resume with
\code{"finetune"}, but you lose the optimiser's moments and the scheduler's position. Consider
whether that is what you want.}
\confusion{Why is the validation loss averaged over patch sizes?}{Because the thermal models train
on several sizes. Indradhanu has one, so the average is over one value.}
\end{confusions}

\begin{checkyourself}
\item\label{ex:trainloop:1} By hand: run three steps of gradient descent on $L(w) = (w - 3)^2$
from $w_0 = 0$ with $\eta = 0.25$. What factor multiplies the distance to 3 at each step? For
which $\eta$ does one step land exactly on 3?
\item\label{ex:trainloop:2} By hand: a parameter's gradient is the constant $g = 2$ at every
step. Compute Adam's $m_t$, $v_t$, $\hat{m}_t$, $\hat{v}_t$ for $t = 1, 2$ with the default
$\beta$s, and the size of the step with $\eta = 10^{-3}$. What would the step have been at $t =
1$ without bias correction?
\item\label{ex:trainloop:3} By hand: compute the v0.2.0 learning rate at epoch indices 25 and
175 ($\cos(\pi/8) = 0.923880$). Which of the two epochs is the rate changing faster in, and why?
\item\label{ex:trainloop:4} Clip $\vect{g} = (1, 2, 2)$ to norm 1.5, and then to norm 5. Then
explain in two sentences why the repository's training would or would not be changed by adding
gradient clipping at norm 1.0, given what Adam does to a gradient's scale.
\item\label{ex:trainloop:5} List the row and column patch coordinates that
\code{PatchPlanGenerator} produces for a 400-row, 300-column scene with size 128 and stride 64,
and the number of patches. Then give a height for which it emits a duplicate row.
\item\label{ex:trainloop:6} Explain why a reconstruction error of 0.0004 at reflectance 0.2
cannot be measured in bfloat16, and what the SpectralEarth trainer does about it.
\end{checkyourself}

\nextchapter{ch:pretrain}
\par\endgroup
